% !TEX program = xelatex
\documentclass[11pt]{ctexart}

% ==================== packages ====================
\usepackage[margin=1in]{geometry}
\usepackage{amsmath,amssymb,amsthm}
\usepackage{graphicx}
\usepackage{booktabs}
\usepackage{multirow}
\usepackage{array}
\usepackage{xcolor}
\usepackage[numbers,sort&compress]{natbib}
\usepackage[colorlinks=true,linkcolor=blue,citecolor=blue,urlcolor=blue]{hyperref}
\usepackage{caption}
\usepackage{etoolbox}
\captionsetup{font=small,skip=4pt}
\setlength{\textfloatsep}{8pt plus 2pt minus 2pt}
\setlength{\floatsep}{6pt plus 2pt minus 2pt}
\setlength{\intextsep}{8pt plus 2pt minus 2pt}
\setcounter{topnumber}{5}
\setcounter{bottomnumber}{5}
\setcounter{totalnumber}{10}
\renewcommand{\topfraction}{0.95}
\renewcommand{\bottomfraction}{0.85}
\renewcommand{\textfraction}{0.05}
\renewcommand{\floatpagefraction}{0.80}
\renewcommand{\abstractname}{\Large 摘要}
\AtBeginDocument{\patchcmd{\abstract}{\vspace{-.5em}}{\vspace{0.8em}}{}{}}
\newcommand{\tablebodysetup}{\footnotesize\setlength{\tabcolsep}{2pt}\renewcommand{\arraystretch}{1.12}}

\newcommand{\raw}{\textsc{raw}}
\renewcommand{\cal}{\textsc{cal}}
\newcommand{\best}[1]{\textbf{#1}}
\newcommand{\E}{\mathbb{E}}
\newcommand{\R}{\mathbb{R}}
\newcommand{\vt}{v_t}
\newcommand{\nfe}{\mathrm{NFE}}

\title{\bfseries 区制切换市场的动作条件流匹配世界模型}

\author{谢若涵\quad 张心一}
\date{课程：基于深度学习的内容生成原理与方法}

\begin{document}
\maketitle

\begin{abstract}
本文研究一类动作条件金融世界模型：外部智能体通过选择正常、高波动、崩盘三类离散市场区制来驱动模拟器，而模型需在不使用教师强制的情况下自回归推演一整个交易年。与仅追求收益直方图逼真的金融生成模型不同，本文关注可控、自洽与定价保真三项要求能否同时成立。为此，本文以三区制马尔可夫切换 Heston 过程构造可控合成市场，并以独立 $100$k 条蒙特卡洛路径计算得到的高精度定价基准评测生成路径上的衍生品价格。评测上，本文不只比较单一价格误差，而是在自由 rollout 下并行报告欧式期权定价 RMSE、亚式期权定价 RMSE、完整价值度曲面稳健性、路径 Wasserstein 距离、风格化事实，以及未校准与矩校准两种定价口径。

方法上，本文沿一条递进路线展开。首先，在经典计量与 GAN 基线之上，本文以两阶段转移流匹配把单步转移分解为方差与收益两个条件流匹配阶段，将未校准定价 RMSE 降至 $0.475$，并通过调度采样考察曝光偏差下的路径统计权衡。这一阶段的结果进一步引出 joint-FM 联合转移流匹配：将下一方差与下一收益作为联合变量一次性建模，使模型直接表达 Heston 中由相关布朗运动产生的方差--收益耦合，并在 Euler $\nfe{=}120$ 下取得未校准定价 RMSE $0.094$。为获得可部署的一步采样器，本文比较 CD、Mean-Flow 与 flow-map 三类一步蒸馏，并提出 on-policy 端点蒸馏：令一步 flow-map 学生在其自身 rollout 访问到的状态分布上匹配冻结 teacher 的端点，将一步模型的 RMSE 由 $0.179$ 降至 $0.101$，同时改善终端收益分布、杠杆相关与峰度。进一步在 $17\times3$ 个欧式价格点和同网格亚式期权上重新选择模型检查点后，最稳健的一步 flow-map 学生体现为高质量 teacher 端点监督、短回放窗口与多轮 on-policy 修正的组合，取得欧式 RMSE $0.0917$ 与亚式 RMSE $0.0525$。

本文进一步将 GARCH-$t$、Quant-GAN、移动块自助法、SIGMA 路径分布微调、可微定价微调，以及 CD、Mean-Flow、flow-map 三类一步蒸馏纳入统一评测，形成从经典基线、两阶段流匹配、联合流匹配到一步部署模型的完整对照。实验表明，联合转移建模与 on-policy 端点蒸馏能够在定价精度、路径统计保真和 NFE1 低成本部署之间取得最稳健的平衡。项目代码与实验产物见 \url{https://github.com/ValkyrieNox/finflow-world-model}。
\end{abstract}

% =====================================================================
\section{引言}
\label{sec:intro}

金融时间序列生成模型被广泛用于压力测试交易策略、扩充稀缺数据，以及在历史中欠采样的情景下为衍生品定价。与图像或文本不同，金融生成器只有在同时复现一整套风格化事实——重尾收益、波动率聚集、杠杆效应——并使其样本路径上的衍生品价格与真实数据生成过程一致时，才具有实用价值~\citep{cont2001empirical}。这两项要求并不等价：一个模型可能边际直方图逼真却为期权定错价格，反之亦然。本文将由独立随机种子生成的 $100$k 条蒙特卡洛路径计算得到的参考期权价格称为高精度蒙特卡洛定价基准，下文简称 MC 定价基准。本文要回答的问题是：如何训练一个既可被外部动作控制、又能在长程自由推演中保持自洽、且生成路径上的期权价格逼近 MC 定价基准的金融世界模型。

\subsection{世界模型视角}
本文采用世界模型的视角~\citep{ha2018world,alonso2024diamond}：生成器是一个动作条件的市场模拟器，外部智能体在每一步通过选择正常、高波动、崩盘三类离散区制来驱动它。本文的状态定义为
\[
s_t=(\log v_t,\ r_{t-1}),
\]
其中 $v_t$ 为 Heston 瞬时方差，$r_{t-1}$ 为上一期对数收益，动作 $a_t$ 为离散市场区制。模型学习单步转移核
\[
p_\theta(\log v_{t+1},\,r_t\mid \log v_t,\,r_{t-1},\,a_t),
\]
并在给定或采样的区制序列下，将该转移核自回归复合 $252$ 次以模拟一个完整交易年。这一框架带来两点贯穿全文的影响。其一，模型在自由 rollout 下评测：自回归展开整年而不回填真实状态，即不使用教师强制。这远比一步预测困难，并暴露曝光偏差，即教师强制训练与自由生成之间的训练与测试错配~\citep{bengio2015scheduled}。其二，可控性要求模型结构必须接受动作通道，并在被强制进入罕见区制时保持稳定。

\subsection{两阶段及联合转移流匹配}
\label{sec:progressive}

\paragraph{两阶段转移流匹配}
相对 GARCH 与 Quant-GAN 等基线，第一步改进是用流匹配学习单步转移。Heston 中方差先于收益生成，存在天然的因果结构，因此最直接的做法是将联合转移分解为先后两个条件流匹配阶段：先采样下一方差，再以其为条件采样下一收益。配合用以缓解曝光偏差的调度采样，两阶段流匹配已明显优于上述基线。

\paragraph{联合转移流匹配}
然而两阶段分解人为割裂了方差与收益由相关布朗运动产生的强耦合：前一阶段的偏差会作为条件进入后一阶段，并在一整年的自回归推演中逐步累积。本文据此从架构层面重新设计转移模型——把下一方差与下一收益作为一个联合变量同时建模，使单步采样一次性给出二者，从结构上保留其相关性。第~\ref{sec:results} 节的实验表明，该联合转移流匹配的定价质量大幅超越两阶段分解。这一“分解 $\to$ 联合”的递进，是本文方法论的核心线索。

\subsection{面向部署的一步蒸馏}
一个市场模拟器只有在长程推演成本足够低时才具备实用性。高质量的转移采样通常需在每一步反复调用网络若干次，若对一整年的每个交易日都如此，代价过高。为便于部署，本文目标是得到只需单次网络调用即可给出下一状态的采样器，做法是借助知识蒸馏，把多步转移模型蒸馏为一个一步模型。下文沿用蒸馏的惯用称谓：提供监督信号的高质量多步模型称作 teacher，被训练去模仿其输出的轻量一步模型称作 student。核心困难在于训练与部署之间的分布错配：若一步模型仅在取自真实数据的条件上训练，部署时它面对的却是自身先前各步生成的状态，这一错配会在长程推演中被放大。本文采用的策略是使一步模型在其自身推演实际访问到的状态上进行蒸馏，并以 teacher 在这些状态上的端点作为学习目标——这一 on-policy 端点修正是本文在蒸馏侧的主要创新。

\subsection{本文贡献}
\begin{itemize}\itemsep2pt
\item 本文复现 Quant-GAN 的核心结构基线，并与 GARCH(1,1)-$t$、移动块自助法等经典参照一起纳入统一自由 rollout 评测协议；在同一 Heston 区制市场、同一定价与路径统计口径下，本文方法显著优于这些基线，为后续改进提供清晰参照。
\item 本文以两阶段转移流匹配及调度采样作为起点，进一步发展出 joint-FM 联合转移流匹配 teacher，将下一方差与下一收益作为联合端点同时建模，使方法路线从可解释的分阶段建模推进到更高保真的联合转移建模。该设计直接对应 Heston 模型中方差冲击与价格冲击的相关结构，在自由 rollout 下显著提升长期路径与定价一致性，并取得 \raw{} 定价 RMSE $0.094$。
\item 本文系统考察多种训练侧改进，包括调度采样、严格正常评分的路径分布微调 SIGMA，以及可微定价微调，形成从单步监督、路径分布约束到定价目标优化的训练方法对照，并据此明确金融世界模型需要在状态分布、路径统计与衍生品定价之间保持一致。
\item 本文比较 CD、Mean-Flow 与 flow-map 三类一步蒸馏方法，并提出 on-policy 端点蒸馏，令一步学生在其自身 rollout 状态分布上匹配冻结 teacher 的端点；该方法将一步学生由 $0.179$ 推至 $0.101$，并在保留路径结构的同时实现 NFE1 部署。
\item 本文构建更完整的评测体系，在原有 $15$ 点欧式期权协议之外加入亚式期权、$17\times3$ 完整价值度曲面、Wasserstein 距离、风格化事实以及 \raw{}/\cal{} 并行报告，使最终模型选择同时反映欧式定价、路径依赖 payoff、尾部价值度区域与路径统计保真。
\end{itemize}

% =====================================================================
\section{相关工作}
\label{sec:related}

\paragraph{金融时间序列生成模型}
Quant-GAN~\citep{wiese2020quant} 以时序卷积生成器与判别器建模金融路径；本文复现其 TCN-GAN 核心结构，并在本实验实现中结合 Lambert-W 变换~\citep{goerg2015lambert} 与 WGAN-GP~\citep{gulrajani2017improved} 稳定重尾收益上的对抗训练。Sig-WGAN 及其条件变体~\citep{ni2021sigwgan,ni2020conditional} 以基于签名的 Wasserstein-1 距离替换学习到的判别器，利用签名~\citep{lyons2007differential} 将路径泛函线性化的性质。TimeGAN~\citep{yoon2019timegan} 在学习到的嵌入空间中结合对抗与监督损失。经典计量基线——GARCH 及其非对称 GJR 变体配以 Student-$t$ 创新~\citep{bollerslev1986garch,glosten1993relation,bollerslev1987conditionally}——在波动率聚集上仍具有竞争力，是深度模型需要比较的参照。

\paragraph{流匹配与世界模型}
流匹配~\citep{lipman2023flow}、修正流~\citep{liu2023rectified} 与随机插值~\citep{albergo2023stochastic} 学习将噪声输运到数据的连续时间速度场；与扩散模型中的分数建模~\citep{ho2020denoising,karras2022elucidating} 和 GAN 的对抗训练相比，这一路线以速度场回归给出更直接的训练目标，并便于后续少步采样与蒸馏。世界模型~\citep{ha2018world,alonso2024diamond} 学习环境的动作条件模拟器；本文将这一视角引入量化金融，把市场区制视为动作。

\paragraph{严格正常评分规则与路径分布损失}
若一个评分规则的期望仅在真实分布处唯一取得最小值，则称其为严格正常的~\citep{gneiting2007strictly}。能量分数以及由特征核诱导的 kernel score/MMD 目标可构成严格正常评分规则~\citep{gretton2012kernel}，并已被用于训练概率预报器~\citep{pacchiardi2024probabilistic} 以及无需判别器的神经 SDE~\citep{issa2023nonadversarial}。本文将签名核 MMD~\citep{salvi2021signature}、Sig-$W_1$ 期望签名项与能量分数等路径级分数组合为一个非对抗的路径分布微调目标。

\paragraph{少步蒸馏}
一致性模型~\citep{song2023consistency,song2024improved} 学习沿概率流 ODE 不变的一步映射。Mean-Flow~\citep{geng2025meanflow} 学习用于一步采样的时间平均速度。Flow-map 自蒸馏~\citep{boffi2025flowmaps} 通过拉格朗日目标学习两时刻输运映射，无需空间导数，并支持 $K$ 步组合以换取质量与算力的可调权衡。本工作在自由 rollout 的金融模拟下，将一致性蒸馏、Mean-Flow 平均速度参数化与 flow-map 作为同一 teacher 的一步学生基线进行比较，并提出 on-policy 端点修正以弥合训练与部署条件分布错配。

% =====================================================================
\section{问题设定与评测协议}
\label{sec:setup}

\subsection{三区制马尔可夫切换 Heston 数据}
真实市场服从 Heston 随机波动率模型~\citep{heston1993}，其参数按马尔可夫链在三个潜在区制间切换。在区制 $a_t\in\{\text{正常},\text{高波动},\text{崩盘}\}$ 下，
\begin{align}
dS_t &= \mu_{a_t} S_t\,dt + \sqrt{v_t}\,S_t\,dW^S_t, \\
dv_t &= \kappa_{a_t}(\theta_{a_t}-v_t)\,dt
      + \xi_{a_t}\sqrt{v_t}\,dW^v_t,\qquad
d\langle W^S,W^v\rangle_t = \rho_{a_t}\,dt .
\end{align}
三区制参数设定为
\[
\begin{array}{lccccc}
\toprule
\text{区制} & \kappa & \theta & \xi & \rho & \mu \\
\midrule
\text{正常} & 2 & 0.04 & 0.3 & -0.7 & 0.05 \\
\text{高波动} & 3 & 0.09 & 0.5 & -0.7 & 0.05 \\
\text{崩盘} & 4 & 0.16 & 0.8 & -0.7 & -0.2 \\
\bottomrule
\end{array}
\]
区制由行随机转移矩阵演化，对角元素 $0.992,0.95,0.90$ 占优以保证持续性。本文采用 Andersen QE 格式~\citep{andersen2008simple} 离散方差，时间步 $dt=1/252$，每条路径包含 $252$ 个交易日，初始价格 $S_0=100$、初始方差 $v_0=0.04$。训练、验证、测试路径分别为 $50$k、$5$k、$10$k；并以独立随机种子生成 $100$k 条 MC 路径构造 MC 定价基准。

\subsection{世界模型的单步转移分解}
本文自回归地建模单步转移，并在全文中比较两种分解。第一种是两阶段分解，按 Heston 的因果结构先后采样方差与收益：
\begin{equation}
p_\theta(\log v_{t+1},r_t\mid\cdot)
= \underbrace{p^{\mathrm{vol}}_\theta(\log v_{t+1}\mid \cdot)}_{\text{方差阶段}}\;
  \underbrace{p^{\mathrm{ret}}_\theta(r_t\mid \log v_{t+1},\cdot)}_{\text{收益阶段}},
\end{equation}
其中 $r_{-1}=0$。第二种是联合分解，将下一方差与下一收益作为一个二维联合变量一次性建模：
\begin{equation}
p_\theta\big(\log v_{t+1},\,r_t \mid \log v_t,\,r_{t-1},\,a_t\big).
\end{equation}
两种分解的具体训练算法见第~\ref{sec:method} 节。一次完整模拟由自由 rollout 产生：采样下一状态、追加，并在外部给定的区制序列下重复 $252$ 步。

\subsection{自由 rollout 与基准定价}
\label{sec:eval}
所有模型均在自由 rollout 下评测。模型自 $(S_0,v_0)$ 出发，自回归生成 $252$ 个交易日，过程中不读取测试集真实状态。评测协议由三组指标构成。第一组是衍生品定价指标：基础协议计算生成路径上欧式看涨期权价格相对 MC 定价基准的 RMSE 与 MAPE，价值度 moneyness 取 $\{0.85,0.90,0.95,1.00,1.05\}$，到期取 $\{0.25,0.5,1.0\}$，共 $15$ 个价格点；同一网格上还计算算术平均亚式看涨期权 RMSE/MAPE，用于考察路径平均 payoff。为检验尾部价值度区域的稳健性，最终检查点选择进一步使用完整价值度曲面协议，将 moneyness 扩展为 $17$ 个点，覆盖 $0.50$ 到 $2.00$，并在同一三个到期上同时报告欧式与亚式 RMSE。第二组是路径分布指标，包括终端收益 Wasserstein $W_1(R_T)$、终端绝对收益 Wasserstein $W_1(|R_T|)$、逐日收益边际 Wasserstein 以及签名 Wasserstein Sig-$W_1$。第三组是风格化事实指标，包括绝对收益自相关 L1 即 Abs-ACF L1、杠杆相关 L1 即 Lev L1，以及峰度。所有期权价格均由蒙特卡洛在生成路径上估计，并与同一 MC 定价基准逐点比较。

\subsection{原始与校准两种报告口径}
\label{sec:rawcal}
本文报告两种定价口径。\raw{} 指原始未校准口径，即直接使用生成器输出的收益与价格路径计算期权价格；它检验的是模型自身学到的转移动力学是否已经具备定价能力，因此是本文的主要口径。本文的主要 teacher、一步蒸馏和 on-policy 修正均首先在 \raw{} 下比较；joint-FM 与 on-policy flow-map 也正是在未经后处理时取得核心结果。\cal{} 指矩校准口径，即在生成器输出之上施加一次仿射后处理：将池化生成收益标准化后，按训练数据收益的均值与标准差重新缩放，再重建价格路径并定价。该处理与 Quant-GAN 所用的矩校准一致，对所有生成器统一适用，不使用 MC 定价基准或测试集价格来调参，因此是合理的辅助评测口径。需要注意的是，矩校准注入了收益的一阶与二阶矩，而期权价格对这两阶矩敏感；所以 \cal{} 衡量的是标准后处理后的定价可用性，不能替代 \raw{} 对生成动力学的检验。在基础 $15$ 点欧式协议下，真实测试集相对 MC 定价基准的 \raw{} RMSE 为 $0.1646$、MAPE 为 $0.0112$、峰度为 $4.60$，可作为 $10$k 条测试路径相对 $100$k 基准路径的有限样本 MC 误差参照。该数值用于标定有限样本误差量级，不作为模型性能的理论下界；完整价值度曲面与亚式期权协议分别报告自身的 RMSE。

% =====================================================================
\section{方法}
\label{sec:method}

本文的方法分为两部分：第~\ref{sec:train} 节讨论转移模型的训练，第~\ref{sec:distill} 节讨论面向部署的一步蒸馏。训练部分沿 \S\ref{sec:progressive} 的递进路线展开，先给出两阶段流匹配及路径分布微调、可微定价微调两种事后微调，再给出从架构上解决其瓶颈的联合转移流匹配 teacher。蒸馏部分按方法演进顺序展开，先比较三类一步蒸馏的参数化，再给出 on-policy 端点修正。两部分共同构成一套面向“可控、自洽、定价保真”这一目标的转移建模与一步部署算法。

\subsection{转移模型的训练}
\label{sec:train}

\subsubsection{两阶段转移流匹配与调度采样}
\label{sec:twostage}
利用 Heston 中方差先于收益生成的因果结构，两阶段流匹配把单步转移分解为方差阶段
$p^{\mathrm{vol}}_\theta(\log v_{t+1}\mid \log v_t,a_t)$
与收益阶段
$p^{\mathrm{ret}}_\theta(r_t\mid \log v_{t+1},\log v_t,r_{t-1},a_t)$，
每一阶段用流匹配~\citep{lipman2023flow,liu2023rectified} 训练一个条件速度场：在噪声 $x_0\sim\mathcal N(0,I)$ 与目标 $x_1$ 之间构造线性插值 $x_\tau=(1-\tau)x_0+\tau x_1$，并以 $\|v_\theta(x_\tau,\tau,c)-(x_1-x_0)\|_2^2$ 回归条件期望速度。

由于 teacher 以一步监督训练却需自由部署，收益阶段在训练时看到的 $\log v_{t+1}$ 通常是真实下一方差，而部署时该条件来自上一方差阶段的模型样本，因而会产生条件分布偏移。为缓解这一问题，本文以调度采样~\citep{bengio2015scheduled} 训练收益阶段：以概率 $p$ 用方差 sampler 生成的 $\log v_{t+1}$ 替换真实下一方差，使收益模型在训练时暴露于部署时会遇到的方差条件误差。$p=0.8$ 给出最佳风格化事实平衡，此时峰度为 $4.41$，并在表~\ref{tab:unified} 的基础协议下取得 \cal{} $0.348$。同一两阶段框架下，FM teacher 的 \raw{} 定价 RMSE 为 $0.475$，显著优于 RMSE 为 $1.089$ 的 GARCH-$t$ 与 RMSE 为 $5.4$ 的 Quant-GAN；但其绝对误差仍高，提示瓶颈在分解本身，而非任一阶段的拟合。基于这一结果，本文进一步考察两类微调方法，并最终转向联合转移架构。

\subsubsection{SIGMA 路径分布微调}
\label{sec:sigma}
两阶段 teacher 的逐步监督只约束单步转移，并不直接约束整条路径的分布。为在路径层面对齐生成分布与真实分布，本文提出 SIGMA，即 strictly-proper-scoring path-distribution fine-tuning。在可微的自由 rollout 上，SIGMA 以严格正常评分规则~\citep{gneiting2007strictly} 启发的路径分布损失对 teacher 做后微调。与 Quant-GAN 依赖学习到的对抗判别器不同，SIGMA 完全非对抗：能量分数与特征核 MMD 等目标可作为分布比较的正常评分或二样本距离，因而可直接转化为优化信号，既无需训练判别器，也避免了对抗训练的不稳定。这是本文在训练侧的创新之一。

具体地，SIGMA 在长度 $L$ 的生成路径分布 $\hat P$ 与真实路径分布 $P$ 之间最小化一组可选路径分布项与稳定正则的加权和：
\begin{equation}
\mathcal L_{\mathrm{SIGMA}}
= \lambda_{\mathrm{sk}}\,\mathrm{MMD}^2_{\mathrm{sigker}}(\hat P,P)
+ \lambda_{\mathrm{s1}}\left\|\E_{\hat P}\Phi_{\mathrm{sig}}(X)-\E_{P}\Phi_{\mathrm{sig}}(X)\right\|_2^2
+ \lambda_{\mathrm{e}}\,\mathrm{Energy}(\hat P,P)
+ \lambda_{\mathrm{aux}}\,\mathcal R_{\mathrm{aux}} .
\end{equation}
其中签名核 MMD~\citep{salvi2021signature,issa2023nonadversarial,gretton2012kernel} 在实现中由截断签名特征上的 RBF-MMD 近似，对波动率聚集等高阶时序依赖敏感；期望签名项~\citep{ni2021sigwgan,ni2020conditional} 比较两分布的平均签名特征，线性地刻画路径泛函；能量分数~\citep{pacchiardi2024probabilistic} 比较整条标准化收益路径；$\mathcal R_{\mathrm{aux}}$ 包含收益均值/方差、终端收益、绝对收益累计、峰度与锚定正则，用于稳定训练并防止尺度漂移。整条 rollout 关于网络参数可微，梯度经重参数化采样回传至 teacher。

SIGMA 能单调改善校准后定价，最优 \cal{} 为 $0.180$，表明路径分布损失能够缩小生成分布与真实分布之间的差异；但其 \raw{} 质量与峰度仍不及后文的 joint-FM。这说明事后的路径微调难以弥补两阶段分解在结构上引入的误差，从而促使本文从“微调既有 teacher”转向“重新设计 teacher 架构”。

\subsubsection{可微定价微调}
\label{sec:diffprice}
SIGMA 在路径分布层面对齐，而相对 MC 定价基准的期权价格误差是本文的重要定价指标，因此可以直接对价格目标求导。可微定价微调将生成路径上的欧式期权价格以可微蒙特卡洛实现，并以逐点相对价格误差的平方均值作为损失；实验中主要从 joint flow-map 学生出发，同时保留 flow-map 蒸馏项作为动力学正则。其优势在于直接优化价格相关目标、且价格关于路径可微，梯度信号明确。本文把它作为一种探索方案纳入，并在 \S\ref{sec:res_pricing} 给出具体消融实验。

本文讨论该方法的原因在于，它刻画了低维定价目标的作用边界：期权价格仅是路径分布在 $15$ 个价格点上的低维投影，直接拟合该投影可将定价 RMSE 降至最优 $0.158$，但实验中也会放大绝对收益自相关与杠杆相关、降低峰度。因此本文把可微定价微调作为局部价格敏感性诊断，而最终模型仍需在路径分布、状态分布与定价指标之间联合选择。这一消融为后文的联合建模与 on-policy 蒸馏提供了直接动机。

\subsubsection{联合转移流匹配教师模型}
\label{sec:jointfm}
两阶段分解、路径微调与定价微调的实验共同指向一个架构改进方向：Heston 中方差与收益由相关布朗运动共同驱动，若把二者拆成先后两个阶段，前一阶段的偏差会作为条件进入后一阶段，并在 $252$ 步自回归推演中累积。本文据此从架构层面重新设计 teacher，将下一方差与下一收益作为一个二维联合变量一次性建模。令
\[
x_1 = (\log v_{t+1},\, r_t), \qquad c=(\log v_t,\,r_{t-1},\,a_t),
\]
在噪声 $x_0\sim\mathcal N(0,I)$ 与真实样本 $x_1$ 之间构造修正流插值~\citep{lipman2023flow,liu2023rectified,albergo2023stochastic}
\[
x_\tau=(1-\tau)\,x_0+\tau\,x_1,\qquad \tau\sim U[0,1],
\]
并训练一个联合速度场以匹配条件期望速度
\begin{equation}
\mathcal L_{\mathrm{joint}}
=\E_{\tau,x_0,x_1,c}
\left\|v_\theta(x_\tau,\tau,c)-(x_1-x_0)\right\|_2^2 .
\end{equation}
速度场为带正弦时间嵌入的残差 MLP，并对区制动作做一定比例的 dropout，以便采样时可用无分类器引导来增强可控性。采样时自 $x_0$ 出发积分 ODE $\dot x=v_\theta(x,\tau,c)$，本文最优 teacher 取 EMA epoch 60、Euler $\nfe{=}120$。

联合建模相较两阶段的关键优势在于：方差与收益的耦合，尤其是杠杆相关，由单个速度场一次性给出，从而减少阶段间条件误差在自由 rollout 中的传递。这一架构改动把 \raw{} 定价 RMSE 从两阶段的 $0.475$ 降至 $0.094$，并使峰度等风格化事实接近真实值，无需任何事后矩校准或定价微调。换言之，本文并非把若干已有模块简单组合，而是针对“联合转移”这一任务结构，构造出一个以联合速度场为核心、以动作 dropout 实现可控性的训练算法，这构成本文方法的主要研究价值。

\subsection{一步蒸馏}
\label{sec:distill}
联合 teacher 质量虽高，但 $\nfe{=}120$ 不适合作为部署采样器：一整年 $252$ 步、每步上百次网络调用，代价过高。蒸馏的目标是得到一个一步学生 $F_\phi(z,c)\approx x_1$，对每个交易日仅需一次前向，即 NFE1。本文按方法演进顺序展开：先比较三类一步蒸馏的参数化，确定表现最佳的初始学生；再在该起点上提出 on-policy 端点修正。

\subsubsection{三类一步蒸馏的参数化}
\label{sec:distillfamily}
在同一冻结的 joint-FM teacher 上，本文比较三种学生参数化。CD 一致性蒸馏~\citep{song2023consistency,song2024improved} 学习沿 teacher 概率流 ODE 不变的一步映射；Mean-Flow~\citep{geng2025meanflow} 回归用于一步采样的时间平均速度；flow-map 自蒸馏~\citep{boffi2025flowmaps} 通过拉格朗日目标直接学习从噪声与条件到 teacher 端点的两时刻映射，无需空间导数。三者中，flow-map 直接以 teacher 的联合端点为目标，与本文对 \raw{} 质量的要求最为一致；\S\ref{sec:res_distill} 的实验表明，基础 flow-map 的一步 RMSE 为 $0.179$，优于 RMSE 为 $0.447$ 的 CD 与 RMSE 为 $2.292$ 的 Mean-Flow，因此本文以 flow-map 作为后续修正的初始学生。

尽管基础 flow-map 是三类学生中表现最好的起点，其 $0.179$ 与 teacher 的 $0.094$ 仍有差距。这一差距并非源于一步参数化的表达能力不足——flow-map 已能较好地表示 teacher 的联合端点——而是源于训练分布与部署分布的不一致：基础 flow-map 在取自真实数据的条件上训练，部署时却面对自身 rollout 产生的状态。这一诊断引出下文的 on-policy 端点修正。

\subsubsection{On-policy 教师端点修正}
\label{sec:onpolicy_method}
本文在蒸馏侧的主要创新，是把训练条件从“真实数据状态”替换为“学生自身 rollout 实际访问的状态”，并在这些状态上用冻结 teacher 提供端点监督，从而直接针对上述曝光偏差。算法如下：
\begin{enumerate}\itemsep2pt
\item 令当前一步学生自由 rollout 若干步，收集学生自身产生的 on-policy 条件 $c_{\mathrm{op}}$；
\item 对同一噪声 $z$ 与同一条件 $c_{\mathrm{op}}$，用冻结 joint-FM teacher 以 $\nfe{=}120$ 积分得到端点 $T(z,c_{\mathrm{op}})$，作为该状态下的监督标签；
\item 训练一步学生输出 $F_\phi(z,c_{\mathrm{op}})$ 匹配该端点，并保留基础 flow-map loss 作为正则；
\item 周期性地用更新后的学生重新收集 $c_{\mathrm{op}}$，使训练条件持续跟随部署分布。
\end{enumerate}
目标函数为
\begin{equation}
\mathcal L_{\mathrm{op}}
=\left\|F_\phi(z,c_{\mathrm{op}})-T(z,c_{\mathrm{op}})\right\|_2^2
+\lambda_{\mathrm{fm}}\,\mathcal L_{\mathrm{flowmap}} .
\end{equation}
该模型在部署时仍为一步学生；teacher 只在训练阶段提供端点标签，不进入部署，因而不增加任何推理成本。与 CD 等多轮 progressive distillation 不同，on-policy 端点修正不引入中间 teacher，也不把目标改成低维价格，而是同时在“学生实际访问的状态分布”和“teacher 的高质量端点”两处对齐，从而既弥合曝光偏差、又不破坏路径动力学。

这一训练目标决定的是如何修正学生的状态分布；最终检查点则按第~\ref{sec:eval} 节的统一评测体系选择。基础 $15$ 点欧式协议用于和全部方法横向对齐：teacher $\nfe{=}120$、on-policy 回放长度 $128$、每次 $512$ 条路径、每轮 $30$ 个更新步、训练 $1$ 轮、学习率 $5\times10^{-6}$、flow-map 正则权重 $1$ 的配置把一步学生的定价 RMSE 由 $0.179$ 降至 $0.101$。完整模型选择进一步同时考察终端欧式 payoff、路径平均亚式 payoff、$0.50$--$2.00$ 的完整价值度曲面以及路径分布距离；在这一更完整的坐标下，最终采用同一 $120$ 步 teacher、回放长度 $64$、每轮 $30$ 个更新步并训练 $4$ 轮的配置，取得欧式 RMSE $0.0917$ 与亚式 RMSE $0.0525$。这样处理与 on-policy 方法本身是一致的：既然训练要对齐学生 rollout 访问到的状态分布，评测也应同时覆盖终端价格、路径依赖 payoff、尾部价值度区域与路径统计，而不是只依赖少量平值附近欧式价格点。

% =====================================================================
\section{实验}
\label{sec:results}

\paragraph{基线}
除上述学习型方法外，本文与三个标准参照对比：其一为 Quant-GAN~\citep{wiese2020quant} 核心结构基线，以 TCN 生成器/判别器复现金融时间序列 GAN，并在本实验实现中结合 Lambert-W 变换与 WGAN-GP 稳定化训练，同时报告其最优与最后检查点；其二为 GARCH(1,1)-$t$~\citep{bollerslev1986garch,bollerslev1987conditionally}，对池化训练收益做极大似然拟合，估计量为 $\hat\alpha{=}0.082$、$\hat\beta{=}0.911$、$\hat\nu{=}6.1$；其三为移动块自助法~\citep{kunsch1989jackknife,politis1994stationary}，对真实训练收益按块长 $21$ 重采样。后两者是经典计量参照，无需 GPU；二者均非动作条件，无法被驱动进入指定区制，而这正是世界模型新增的可控性。

\subsection{联合转移建模与一步部署性能}
\label{sec:res_main}
表~\ref{tab:main_joint} 给出核心结果，并印证 \S\ref{sec:progressive} 的递进路线。沿“两阶段 $\to$ 联合”这一主线，定价质量出现显著改善：表~\ref{tab:unified} 中作为第一阶优化的两阶段流匹配 teacher 为 $0.475$，而联合的 joint-FM teacher 降至 $0.094$。这一改善主要来自联合速度场对方差与收益的共同建模，使 teacher 更直接地表达 Heston 的相关冲击结构。joint-FM teacher 因此构成后续一步蒸馏要逼近的多步采样参照；其峰度 $4.371$ 也已接近真实值 $4.60$，且全程未用任何矩校准或定价微调，说明该质量来自学到的动力学本身。

\begin{table}[htbp]
\centering
\caption{核心结果：joint-FM teacher 与一步学生，均为未校准 \raw{} 评测。基础 $15$ 点欧式协议下，真实测试集相对 MC 定价基准的 RMSE $0.1646$ 为有限样本误差参照；teacher 使用 Euler NFE120；所有学生均为 NFE1 一步部署。每列最优值以粗体标记。}
\label{tab:main_joint}
\tablebodysetup
\begin{tabular}{l c c c c c c c}
\toprule
模型 & RMSE $\downarrow$ & MAPE $\downarrow$ & $W_1(R_T)$ $\downarrow$ & 边际 $W_1$ $\downarrow$ & Abs-ACF L1 $\downarrow$ & Lev L1 $\downarrow$ & 峰度 $\to4.60$ \\
\midrule
真实测试集 vs MC 基准 & 0.1646 & 0.0112 & -- & -- & -- & -- & 4.603 \\
Joint-FM teacher & \best{0.0942} & \best{0.0095} & 0.0061 & 0.00030 & 0.0043 & 0.00086 & \best{4.371} \\
Flow-map & 0.1789 & 0.0145 & 0.0131 & 0.00051 & \best{0.0042} & 0.00212 & 5.974 \\
Pricing-aware flow-map & 0.1577 & 0.0174 & 0.0236 & 0.00080 & 0.0430 & 0.00892 & 3.350 \\
On-policy flow-map & 0.1010 & 0.0110 & \best{0.0052} & \best{0.00028} & 0.0047 & \best{0.00086} & 4.356 \\
\bottomrule
\end{tabular}
\end{table}

蒸馏结果也显示出同样的趋势。基础 flow-map 一步模型的 RMSE 为 $0.179$，继承了 teacher 的部分能力但仍有可见差距，且其峰度偏高至 $5.974$、终端 $W_1$ 偏大至 $0.0131$，表明曝光偏差会在长 rollout 中累积。on-policy 端点修正同时改善价格与路径分布：RMSE 由 $0.179$ 降至 $0.101$，终端收益 Wasserstein 由 $0.0131$ 降至 $0.0052$，杠杆相关由 $0.00212$ 降至 $0.00086$，峰度由 $5.974$ 调至 $4.356$，接近真实值 $4.60$。这表明把训练条件对齐到学生实际访问的状态分布，能够在多个相互独立的维度上同时改善结果，而非仅降低单一指标。该学生在 $W_1(R_T)$ 上与 teacher 的评测样本相当，应理解为有限样本波动，而非学生整体优于 teacher。pricing-aware flow-map 则形成另一种权衡：它把局部欧式 RMSE 降至 $0.1577$，但 Abs-ACF L1 提高到 $0.0430$、峰度降低到 $3.350$，因此更适合作为定价敏感性消融，而不是最终部署学生。

\subsection{三类一步蒸馏方法比较}
\label{sec:res_distill}
在 joint-FM teacher 上，本文比较 CD、Mean-Flow 与 flow-map 三类一步蒸馏，结果见表~\ref{tab:distill_joint}。该表在基础 $15$ 点欧式协议下同时报告欧式定价、同网格亚式定价、路径分布与风格化事实指标。三者呈现清晰分工：CD 在终端收益 $W_1$ 与逐日边际 $W_1$ 上略低，Mean-Flow 的峰度最接近真实值 $4.60$；但二者的价格与路径依赖 payoff 误差明显更大。CD 的欧式 RMSE 为 $0.447$、亚式 RMSE 为 $0.274$，Mean-Flow 的欧式 RMSE 为 $2.292$、亚式 RMSE 为 $0.970$，均落后于基础 flow-map 的欧式 RMSE $0.179$ 与亚式 RMSE $0.052$。

\begin{table}[htbp]
\centering
\caption{Joint-FM teacher 上的三种一步蒸馏。除 Asian RMSE 为算术平均亚式看涨 payoff 外，其余各列均为基础 $15$ 点欧式协议下的结果：RMSE 与 MAPE 为欧式期权定价指标，右侧路径统计由同一自由 rollout 计算。每列最优值以粗体标记；峰度列按最接近真实值 $4.60$ 标记。}
\label{tab:distill_joint}
\tablebodysetup
\begin{tabular}{l c c c c c c c c}
\toprule
一步学生 & RMSE $\downarrow$ & MAPE $\downarrow$ & Asian RMSE $\downarrow$ & $W_1(R_T)$ $\downarrow$ & 边际 $W_1$ $\downarrow$ & Abs-ACF L1 $\downarrow$ & Lev L1 $\downarrow$ & 峰度 $\to4.60$ \\
\midrule
CD & 0.447 & 0.0361 & 0.274 & \best{0.0113} & \best{0.00037} & 0.0072 & 0.00235 & 4.028 \\
Mean-Flow & 2.292 & 0.1929 & 0.970 & 0.0406 & 0.00164 & 0.0046 & 0.00569 & \best{4.624} \\
Flow-map & \best{0.179} & \best{0.0145} & \best{0.052} & 0.0131 & 0.00051 & \best{0.0042} & \best{0.00212} & 5.974 \\
\bottomrule
\end{tabular}
\end{table}

因此，本文选择 flow-map 并不是因为它在每一个低维统计上都逐列最优，而是因为它在后续一步部署最关键的指标组合上最稳健：它同时取得最低欧式定价误差、最低亚式定价误差、最低 Abs-ACF L1 与最低杠杆相关 L1。终端或边际 Wasserstein 是重要的分布距离，但只刻画收益分布的低维投影；峰度也只是单一矩统计。相比之下，欧式价格、路径平均的亚式 payoff、波动率聚集和杠杆相关共同约束了路径在时间维度上的金融结构。表~\ref{tab:distill_joint} 因而支持将 flow-map 作为后续 pricing-aware 与 on-policy 修正的初始学生。

flow-map 更适合作为本文的一步学生起点，原因在于优化目标与 joint-FM teacher 的端点结构一致：teacher 学的是把噪声输运到联合端点 $x_1=(\log v_{t+1},r_t)$，而 flow-map 直接学习两时刻输运映射并在一步采样时输出该端点；CD 约束的是沿 ODE 的不变性、Mean-Flow 回归的是时间平均速度，二者都要在长复合 rollout 中把局部一致性累积成全局端点精度，误差更易放大。这一比较说明本文后续所有 pricing-aware 与 on-policy 修正均建立在 flow-map 之上，而非 CD 或 Mean-Flow。

\subsection{可微定价微调的消融}
\label{sec:res_pricing}
表~\ref{tab:pricing} 展示 pricing-aware flow-map 可微定价微调的三个配置。温和配置 $w{=}10,\lambda_{\mathrm{dyn}}{=}0.2$ 将欧式 RMSE 从基础 flow-map 的 $0.179$ 降至 $0.158$，说明可微蒙特卡洛定价确实能把梯度信号传回生成器，并改善其在训练价格网格上的局部欧式定价。与此同时，该配置的亚式 RMSE 从 $0.052$ 升至 $0.092$，Abs-ACF L1 由 $0.0042$ 提高至 $0.0430$，杠杆相关由 $0.0021$ 提高至 $0.0089$；结合表~\ref{tab:main_joint} 的峰度指标可见，其峰度也由 $5.97$ 降至 $3.35$。

\begin{table}[htbp]
\centering
\caption{可微定价微调消融。除 Asian RMSE 为算术平均亚式看涨 payoff 外，其余各列均为基础 $15$ 点欧式协议下的结果：RMSE 与 MAPE 为欧式期权定价指标，右侧路径统计由同一自由 rollout 计算。温和定价微调可降低局部欧式 RMSE，亚式与路径统计列用于检验该收益是否同步迁移到路径层面。}
\label{tab:pricing}
\tablebodysetup
\begin{tabular}{l c c c c c c c}
\toprule
配置 & RMSE $\downarrow$ & MAPE $\downarrow$ & Asian RMSE $\downarrow$ & $W_1(R_T)$ $\downarrow$ & 边际 $W_1$ $\downarrow$ & Abs-ACF L1 $\downarrow$ & Lev L1 $\downarrow$ \\
\midrule
基础 flow-map & 0.179 & \best{0.0145} & \best{0.052} & \best{0.0131} & \best{0.00051} & \best{0.0042} & \best{0.0021} \\
$w{=}10,\lambda_{\mathrm{dyn}}{=}0.2$ & \best{0.158} & 0.0174 & 0.092 & 0.0236 & 0.00080 & 0.0430 & 0.0089 \\
$w{=}10,\lambda_{\mathrm{dyn}}{=}1.0$ & 0.365 & 0.0453 & 0.189 & 0.0316 & 0.00178 & 0.0368 & 0.0068 \\
$w{=}30,\lambda_{\mathrm{dyn}}{=}0.1$ & 0.747 & 0.0730 & 0.385 & 0.0320 & 0.00150 & 0.0435 & 0.0101 \\
\bottomrule
\end{tabular}
\end{table}

因此，表~\ref{tab:pricing} 的核心信息不是简单否定定价微调，而是刻画它的适用边界：当训练目标只覆盖 $15$ 个欧式价格点时，模型能获得局部欧式定价收益，但该收益未必同步迁移到路径平均 payoff 与路径统计。更强的定价权重或更弱的动力学正则也没有继续带来收益，$\lambda_{\mathrm{dyn}}{=}1.0$ 与 $w{=}30$ 的欧式 RMSE 分别升至 $0.365$ 和 $0.747$，亚式 RMSE 也升至 $0.189$ 和 $0.385$。这组消融支持本文后续选择：定价目标可作为有价值的诊断和微调信号，但最终一步学生更适合通过状态分布上的 teacher 端点监督来保持整体路径动力学。

\subsection{On-policy 端点修正的消融分析}
\label{sec:res_onpolicy}
表~\ref{tab:onpolicy_ablation} 在最终模型选择坐标下扫描 on-policy 端点修正检查点。该坐标同时包含 $17\times3$ 欧式完整价值度曲面、同网格亚式期权，以及边际、终端、绝对终端和签名 Wasserstein 距离；moneyness 覆盖 $0.50$ 到 $2.00$，因而将平值附近、深度实值和深度虚值区域放在同一评测面上。表中配置名依次标出 teacher 采样步数、on-policy 回放长度、每轮更新步数和训练轮数。

\begin{table}[htbp]
\centering
\caption{On-policy flow-map 在 $17\times3$ 完整价值度曲面与亚式期权上的消融。所有配置均为一步模型；欧式与亚式 RMSE 均相对同一 $100$k MC 定价基准计算。}
\label{tab:onpolicy_ablation}
\tablebodysetup
\begin{tabular}{l c c c c c c}
\toprule
配置 & 欧式 RMSE $\downarrow$ & Asian RMSE $\downarrow$ & 边际 $W_1$ $\downarrow$ & $W_1(R_T)$ $\downarrow$ & $W_1(|R_T|)$ $\downarrow$ & Sig-$W_1$ $\downarrow$ \\
\midrule
NFE30/h64/s60/e4 & 0.1322 & 0.0545 & 0.000623 & 0.00810 & 0.1474 & 0.00223 \\
NFE120/h64/s15/e1 & 0.1132 & 0.0694 & 0.000300 & 0.00477 & 0.0450 & 0.00117 \\
NFE120/h64/s30/e1 & 0.1331 & 0.0760 & 0.000294 & 0.00528 & 0.0505 & 0.00127 \\
NFE120/h64/s30/e2 & 0.1071 & 0.0598 & 0.000293 & 0.00470 & 0.0433 & 0.00117 \\
NFE120/h64/s30/e4 & \best{0.0917} & \best{0.0525} & 0.000299 & \best{0.00449} & 0.0437 & \best{0.00111} \\
NFE120/h128/s30/e1 & 0.1416 & 0.0787 & \best{0.000287} & 0.00534 & 0.0438 & 0.00129 \\
NFE120/h252/s30/e1 & 0.1622 & 0.0948 & 0.000301 & 0.00626 & \best{0.0428} & 0.00142 \\
\bottomrule
\end{tabular}
\end{table}

从表中可以看出，使用 $120$ 步 teacher、回放长度 $64$、每轮更新 $30$ 步并训练 $4$ 轮的配置给出了最均衡的终版一步学生：欧式 RMSE 为 $0.0917$、亚式 RMSE 为 $0.0525$，同时取得最低 $W_1(R_T)$ 和最低 Sig-$W_1$。同一 teacher 下，训练轮数从 $1$ 轮增加到 $2$ 轮、再到 $4$ 轮时，欧式与亚式 RMSE 整体改善，说明短回放窗口配合多轮 on-policy 端点修正能够更充分地覆盖学生实际访问的状态分布。其他配置也提供了有用参照：NFE120/h128/s30/e1 在边际 $W_1$ 上最低，NFE120/h252/s30/e1 在绝对终端收益 $W_1$ 上最低，而 NFE30/h64/s60/e4 的较低亚式 RMSE 伴随更大的分布距离。综合定价曲面、路径依赖 payoff 与路径分布三类指标，NFE120/h64/s30/e4 是最稳健的部署选择。

\subsection{与基线及其他方法的统一对比}
\label{sec:res_baselines}
表~\ref{tab:unified} 给出基础 $15$ 点欧式协议下的跨家族汇总，覆盖经典基线、两阶段模型及其微调、joint-FM teacher 和最终 on-policy 一步模型；三类一步蒸馏与 pricing-aware 微调已分别在表~\ref{tab:distill_joint}--\ref{tab:pricing} 展开，这两组实验在同一 $15$ 点网格上的亚式结果也列在相应表中，最终 on-policy 检查点在 $17\times3$ 完整价值度曲面和亚式期权下的选择见表~\ref{tab:onpolicy_ablation}。

\begin{table}[htbp]
\centering
\caption{统一对比：所有方法在 \raw{} 与 \cal{} 定价 RMSE 及峰度上的表现。\raw{} 为未校准、\cal{} 为矩校准后的期权定价 RMSE，数值越低越好；基础 $15$ 点欧式协议下，$0.1646$ 为真实测试集相对 MC 定价基准的有限样本误差参照。峰度目标 $4.60$。$^{\ddagger}$ 非参数化重采样参照：重放真实收益，不可控亦不能外推，是参照而非竞争生成器。}
\label{tab:unified}
\tablebodysetup
\begin{tabular}{l l c c c}
\toprule
& 模型 & \raw{} $\downarrow$ & \cal{} $\downarrow$ & 峰度 $\to4.60$ \\
\midrule
\multirow{2}{*}{\shortstack[l]{参照\\非生成}}
 & 真实测试集 vs MC 基准 & 0.1646 & -- & 4.60 \\
 & 移动块自助法$^{\ddagger}$ 数据重放~\citep{politis1994stationary} & 0.269 & 0.244 & 4.58 \\
\midrule
\multirow{3}{*}{基线}
 & Quant-GAN 最优检查点~\citep{wiese2020quant} & 5.650 & 1.607 & 10.21 \\
 & Quant-GAN 最后检查点~\citep{wiese2020quant} & 5.429 & 0.674 & 6.39 \\
 & GARCH(1,1)-$t$~\citep{bollerslev1986garch} & 1.089 & 0.482 & 19.17 \\
\midrule
\multirow{4}{*}{两阶段与微调}
 & 两阶段 FM teacher & 0.475 & 0.583 & 3.88 \\
 & 两阶段 FM + 调度采样 $p{=}0.8$ & 0.872 & 0.348 & 4.41 \\
 & SIGMA-Pin 路径分布微调最优 \cal{} & 2.595 & 0.180 & 3.63 \\
 & SIGMA-Base $\to$ CD & 1.875 & 0.170 & 3.11 \\
\midrule
\multirow{2}{*}{本文方法}
 & Joint-FM teacher NFE120 & \best{0.094} & \best{0.085} & 4.37 \\
 & On-policy flow-map 一步部署 & 0.101 & 0.092 & 4.36 \\
\bottomrule
\end{tabular}
\end{table}

表~\ref{tab:unified} 可以从四个角度理解。第一，经典参照提供了有意义的评测坐标与对照：移动块自助法的 \raw{} 为 $0.269$、峰度 $4.58$，说明重采样真实收益块可以在无控制、无外推的前提下获得很强的路径统计；GARCH-$t$ 的 \raw{} 为 $1.089$、\cal{} 为 $0.482$，显示无条件波动建模仍有定价价值，但峰度 $19.17$ 明显过高且不支持区制动作。第二，Quant-GAN 两个检查点的 \raw{} 均在 $5.4$--$5.7$，矩校准后分别降至 $1.607$ 与 $0.674$，说明标准后处理能释放其部分定价可用性，也说明检查点选择对定价指标影响很大。第三，两阶段路线本身已经把 \raw{} 推到 $0.475$，调度采样版本则把峰度调整到 $4.41$、\cal{} 降至 $0.348$，体现出训练条件分布修正对校准后定价和风格化事实的价值；SIGMA-Pin 与 SIGMA-Base $\to$ CD 进一步把 \cal{} 推到 $0.180$ 和 $0.170$，展示了路径分布微调加矩校准口径下的竞争力。第四，joint-FM teacher 与 on-policy flow-map 在 \raw{} 下已经达到 $0.094$ 与 $0.101$，校准后仅小幅降至 $0.085$ 与 $0.092$，同时峰度保持在 $4.36$--$4.37$。这说明联合转移建模和 on-policy 端点修正不仅在后处理后可用，也在原始未校准路径上直接具备定价能力。

% =====================================================================
\section{讨论}
\label{sec:discussion}

\subsection{On-policy 修正的作用机制}
三类基础蒸馏方法的比较首先说明，学生参数化本身会强烈影响一步部署质量。CD 在终端收益 $W_1$ 与边际 $W_1$ 上略好，Mean-Flow 的峰度更接近真实值，但二者在欧式与亚式定价上误差明显更大；flow-map 则同时给出最低欧式 RMSE、最低 Asian RMSE、最低 Abs-ACF L1 与最低杠杆相关 L1。因此，flow-map 是最合适的初始一步学生，但它与 NFE120 teacher 之间仍存在可见差距：基础 flow-map 的 RMSE 为 $0.179$、峰度为 $5.974$，说明仅在真实转移条件上训练出的学生，部署到自身自由 rollout 后仍会积累状态分布错配。

On-policy 端点修正正是针对这一错配：它让学生先在自身 rollout 中产生状态，再由冻结 teacher 在这些状态上提供联合端点监督。该方法没有把目标改成低维价格，也没有改变一步部署成本，却在表~\ref{tab:main_joint} 中把 RMSE 从 $0.179$ 降至 $0.101$，并同步改善终端 $W_1$、杠杆相关与峰度。进一步的最终模型选择并非额外补充实验，而是同一评测思想的自然延伸：部署模型需要同时面对欧式终端 payoff、亚式路径平均 payoff、$0.50$--$2.00$ 的价值度区域以及路径分布距离。表~\ref{tab:onpolicy_ablation} 中回放长度 $64$、每轮 $30$ 步、训练 $4$ 轮的配置在这些坐标下取得欧式 RMSE $0.0917$ 与 Asian RMSE $0.0525$，说明状态分布对齐带来的收益能够跨 payoff 类型与价值度区域保持稳定。

\subsection{可微定价微调的作用边界}
可微定价微调的价值在于提供直接、可解释的价格梯度。表~\ref{tab:pricing} 中温和配置把基础 flow-map 的欧式 RMSE 由 $0.179$ 降至 $0.158$，说明生成路径上的欧式期权价格确实可以作为有效的局部优化信号。不过，期权价格只是路径分布的一个投影；当目标只覆盖基础 $15$ 点欧式网格时，局部价格收益并不自动转化为路径依赖 payoff 或时序统计收益。同一配置下 Asian RMSE 从 $0.052$ 升至 $0.092$，Abs-ACF L1 与杠杆 L1 也明显上升；更强定价权重或更弱动力学正则没有继续改善欧式 RMSE。

因此，本文对可微定价微调的结论是边界性的而非否定性的：它适合作为价格敏感性诊断，也可用于检查某个学生是否能响应特定价格目标；但最终部署模型不能只按少量欧式价格点选择。金融世界模型需要同时通过欧式价格、亚式 payoff、路径分布距离与风格化事实的检验，尤其当模型要服务于长程自由 rollout 和下游控制任务时，路径动力学的一致性比单一价格网格更关键。

\subsection{原始与校准的权衡}
\label{sec:tradeoff}
表~\ref{tab:unified} 中的 \raw{}/\cal{} 对照体现了两类有用但不同的能力。\cal{} 衡量模型经标准矩校准后的定价可用性，因此 Quant-GAN、SIGMA 与 CD 在校准后误差明显下降是有实际意义的：它说明这些生成器在经过统一后处理后仍能提供可用的定价样本。与此同时，\raw{} 更直接检验模型本身的自由 rollout 是否已经学到正确的转移动力学与价格尺度。二者并不是互相替代的指标，而是分别回答“后处理后能否定价”和“原始生成过程是否自洽”两个问题。

本文方法的优势在于这两个口径之间的间隙较小。Joint-FM teacher 在 \raw{} 下达到 $0.094$，on-policy flow-map 一步模型达到 $0.101$；矩校准后分别为 $0.085$ 与 $0.092$，改善幅度较小但方向一致。这说明联合转移建模与 on-policy 端点修正已经把主要定价能力嵌入原始路径动力学中，而不是主要依赖事后尺度修正。相反，若某方法 \cal{} 很低但 \raw{} 较高，它仍然可能是有价值的校准后定价工具，只是其应用场景更接近“生成后处理样本”，而不是无需校准的动作条件世界模型。

\subsection{局限与未来工作}
本文仍有两个可扩展之处。第一，衍生品评测目前覆盖欧式看涨与算术平均亚式看涨 payoff，已经分别检验终端价格与路径平均信息，但尚未纳入障碍、回望、篮子等更复杂产品。第二，on-policy 修正依赖高质量 NFE120 teacher 生成端点标签；虽然最终部署端已经是 NFE1，训练阶段仍有进一步压缩 teacher 调用与提升多资产扩展效率的空间。

后续工作可以沿更一般的金融世界模型方向推进：将同一框架迁移到真实股票、期权隐含波动率曲面和限价订单簿数据，并把区制动作扩展为宏观冲击、流动性冲击或策略动作；同时，把正常评分路径损失、定价曲面验证和 on-policy 状态回放纳入同一个训练循环，使 teacher 与学生都在多指标评测体系下共同选择。更长期地，这类动作条件金融世界模型可作为压力测试、策略学习和衍生品风险管理的统一模拟器，而本文的评测协议为这种用途提供了一个可复现起点。

% =====================================================================
\section{结论}
\label{sec:conclusion}

本文提出一个区制切换 Heston 金融世界模型，并围绕“可控、自洽、定价保真”这一目标形成了从基线复现到一步部署的完整方法路线。在 GARCH-$t$、Quant-GAN 与移动块自助法参照之上，两阶段转移流匹配首先把 \raw{} 定价 RMSE 降至 $0.475$；进一步的 joint-FM 联合转移流匹配把下一方差与下一收益作为联合端点建模，使 teacher 在 NFE120 下达到 \raw{} $0.094$。为降低部署成本，本文比较 CD、Mean-Flow 与 flow-map 三类一步蒸馏，选择在欧式定价、亚式定价、波动率聚集与杠杆相关上综合最稳健的 flow-map 作为学生起点，并提出 on-policy 端点蒸馏，使一步模型从 $0.179$ 改善至 $0.101$，在 NFE1 部署下逼近多步 teacher。

实验部分以多维评测而非单一价格误差组织结论：基础 $15$ 点欧式协议用于横向比较全部方法，亚式期权检验路径平均 payoff，$17\times3$ 完整价值度曲面覆盖更宽的尾部定价区域，Wasserstein 距离、Abs-ACF、杠杆相关、峰度以及 \raw{}/\cal{} 并行报告共同刻画路径质量与定价可用性。在这一统一评测体系下，SIGMA 路径分布微调、可微定价微调、CD 与 Mean-Flow 分别展示了路径损失、价格梯度和不同一步蒸馏目标的作用边界；最终的 on-policy flow-map 配置以高质量 teacher 端点、短回放窗口和多轮状态分布修正取得最均衡结果，在完整价值度曲面与亚式期权下达到欧式 RMSE $0.0917$ 与 Asian RMSE $0.0525$。整体结果表明，联合转移建模和学生 rollout 状态分布上的端点蒸馏，是兼顾原始路径动力学、衍生品定价和一步低成本部署的有效方案；本文的评测协议也为后续金融世界模型提供了可复现的比较框架。总体而言，本文提出了一条从高保真联合 teacher 到低成本一步学生的有效训练与蒸馏路线。同时，本文建立了覆盖定价、路径依赖 payoff、尾部价值度区域、路径分布与校准口径的多角度评测范式，为面向压力测试、策略学习和风险管理的金融世界模型提供了更系统的研究起点。

\bibliographystyle{plainnat}
\bibliography{references}

\end{document}
